\subsection{Episode definition}

An episode is defined as
\[
(T, o^\star, q, D, s, g, \phi, \ell, r),
\]
where \(T\) is the simulator task, \(o^\star\) is the target object, \(q\) is the language query, \(D\) is the detection list, \(s\) is the selector or detector, \(g\) is the RGB-D target source, \(\phi\) is the stressor family, \(\ell\) is the stress level, and \(r\) is the seed. The output is a 3D target, diagnostic success, target-error distance, failure type, and auxiliary detector fields when available.

\subsection{Target sources}

The oracle target uses privileged simulator geometry and is used only to estimate feasibility. Box-center depth projects the selected detection center through the depth map. Crop-median depth takes the median valid depth inside the selected crop. Crop-trimmed median depth removes the lowest and highest 10\% of valid crop depths before taking the median, providing a stronger inlier-style RGB-D baseline without adding a new external model dependency. Crop-top-surface is retained as a deliberately simple heuristic in the larger diagnostic sweeps.

\subsection{Selectors and detector bridge}

The metadata query-aware selector uses simulator-provided labels and metadata to select the target detection, and is treated as a controlled metadata-assisted upper-bound comparator rather than deployed perception. It can still fail when query templates and exposed metadata define different targets or when the stressed detection list lacks a target-consistent entry. The first-detection selector selects the first available detection and is used to expose list-order and distractor sensitivity. The GroundingDINO bridge is a learned open-vocabulary detector plug-in; it reports detection availability, IoU summaries when debug fields are present, and no-detection outcomes when no usable detection is returned.

\subsection{Stressor families}

Stressors are grouped into detection-list semantic ambiguity, visual, geometric, and target-displacement families. The same-shape, same-color, and nearby stressors manipulate detection-list ordering and crop ambiguity; they are not claimed to be physically inserted scene objects unless accompanied by rendered scene evidence. Visual and geometric stressors perturb occlusion, depth sparsity, depth noise, and camera pose. Target-displacement stressors perturb the target after generation to audit tolerance to downstream offset. The exact level values are generated from YAML configuration files and reported in Supplementary Table S1.

\subsection{Metrics and confidence intervals}

Diagnostic success is computed from the 3D target-error threshold, with 0.08 m as the default. Threshold sensitivity recomputes success at 0.04 m, 0.06 m, 0.08 m, and 0.10 m from stored target-error values. Oracle gap is
\[
\mathrm{gap}(g)=\mathrm{SR}(\mathrm{oracle})-\mathrm{SR}(g),
\]
where \(\mathrm{SR}\) is diagnostic success rate. Oracle rows use privileged simulator target geometry, but they are still evaluated with the same threshold and stressor protocol as non-oracle rows. Oracle success can therefore be below one under hard target-displacement, projection, visibility, or threshold conditions. Confidence intervals in the main diagnostic tables use seed-block bootstrap aggregation to reduce overconfident episode-level independence assumptions.

\subsection{Result generation and audit}

All numerical tables in this draft are generated from JSON or CSV artifacts. The result-to-claim gate records completed episode counts, duplicate counts, runner exceptions, and whether each intended claim is supported. The closed-loop smoke test is included as an audit because it completed technically but failed the oracle execution gate.
